% =================================================
% Версия для учителя
% =================================================
\worksheettitle
  {Рациональное уравнение с модулем и параметром}
  {Методический сценарий и полный разбор}

\begin{teacherbox}[Назначение комплекта]
  Ученический лист построен как управляемое открытие. Он не сообщает ответ
  непрерывным готовым текстом: ученик формулирует правило нуля дроби,
  раскрывает модуль, сохраняет условия ветвей, проверяет знаменатель,
  собирает ответ на числовой прямой и затем восстанавливает решение целиком.
\end{teacherbox}

\begin{teacherbox}[Входная диагностика]
  Перед основной работой дайте ученикам 3--4 минуты на три коротких задания
  в начале рабочего листа. Они проверяют необходимые действия, но не
  воспроизводят решение основной задачи целиком.

  Ответы: \(x=3\); при \(x\ge0\) имеем \(|x|=x\), при \(x<0\) ---
  \(|x|=-x\); после подстановки получается \(a^2\).
\end{teacherbox}

\begin{teacherbox}[Как использовать результат]
  \begin{itemize}[itemsep=0.25em,topsep=0.25em]
    \item Три верных ответа: сразу переходите к основной задаче и правило
      нуля дроби сформулируйте кратко.
    \item Одна ошибка: адресно разберите соответствующий подготовительный
      приём, не повторяя всю теорию для класса.
    \item Две или три ошибки: сначала выполните вводный раздел фронтально,
      затем вернитесь к основной задаче.
  \end{itemize}
  Диагностика не оценивается отметкой: она нужна для выбора объёма опоры.
\end{teacherbox}

\section{Дидактическая карта}

\begin{center}
  \renewcommand{\arraystretch}{1.38}
  \begin{tabularx}{\textwidth}{
    >{\raggedright\arraybackslash}p{0.26\textwidth}
    >{\raggedright\arraybackslash}X}
    \toprule
    \rowcolor{TableHeader}
    Компонент & Содержание \\
    \midrule
    Главный результат
      & Ученик разделяет исследование на два слоя: число нулей
        числителя и допустимость каждого нуля для исходного уравнения. \\
    \addlinespace[0.25em]
    Предварительные знания
      & Ноль рационального выражения, раскрытие модуля, линейные уравнения,
        подстановка выражения, разложение квадратного трёхчлена. \\
    \addlinespace[0.25em]
    Основной метод
      & Прямой аналитический: две ветви модуля, два кандидата,
        две проверки знаменателя. \\
    \addlinespace[0.25em]
    Дополнительные взгляды
      & Графическое объяснение условия \(a>-2\) и поиск общих корней для
        получения исключаемых параметров. \\
    \addlinespace[0.25em]
    Обязательный маршрут
      & Разделы ученического листа от правила нуля дроби до окончательного
        ответа. Рекомендуемое время: 40--45 минут. \\
    \addlinespace[0.25em]
    Расширение
      & Контрольные значения, графическое объяснение, поиск общих корней,
        самостоятельная реконструкция и перенос метода. Рекомендуется для
        второго урока, домашней работы или сильной группы. \\
    \bottomrule
  \end{tabularx}
\end{center}

\section{Маршрут урока}

\begin{center}
  \renewcommand{\arraystretch}{1.42}
  \begin{tabularx}{\textwidth}{
    >{\centering\arraybackslash}p{0.09\textwidth}
    >{\raggedright\arraybackslash}p{0.24\textwidth}
    >{\raggedright\arraybackslash}p{0.31\textwidth}
    >{\raggedright\arraybackslash}X}
    \toprule
    \rowcolor{TableHeader}
    Время & Этап & Вопрос учителя & Ожидаемый результат \\
    \midrule
    4 мин & Входная диагностика
      & Какие из трёх подготовительных действий уже выполняются уверенно?
      & Учитель выбирает полный или сокращённый вводный разбор. \\
    \addlinespace[0.2em]
    10 мин & Числитель
      & Какие условия нужно сохранить после раскрытия модуля?
      & \(x\ge0\) и \(x<0\); отсюда условия на \(a\). \\
    \addlinespace[0.2em]
    5 мин & Граница
      & Почему при \(a=-2\) две формулы не дают два корня?
      & Значение \(0\) не принадлежит отрицательной ветви. \\
    \addlinespace[0.2em]
    13 мин & Знаменатель
      & Что означает равенство \(D(x_i)=0\)?
      & Кандидат запрещён; исключается соответствующее значение \(a\). \\
    \addlinespace[0.2em]
    5 мин & Сборка ответа
      & Какие условия являются необходимыми и достаточными вместе?
      & \(a>-2\) и \(a\notin\{-1,0,3,8\}\). \\
    \addlinespace[0.2em]
    5 мин & Итоговая проверка
      & Почему полученное условие является необходимым и достаточным?
      & Числитель имеет два различных нуля, и оба допустимы для исходного
        уравнения. \\
    \bottomrule
  \end{tabularx}
\end{center}

\begin{teacherbox}[Граница обязательного маршрута]
  Основной урок заканчивается после записи окончательного ответа в
  ученическом листе. Не требуется в тот же урок проходить контрольную
  таблицу, графический способ, метод общих корней и задания на перенос.
  Эти материалы образуют отдельный блок «Расширение и углубление».
\end{teacherbox}

\begin{teacherbox}[Как использовать расширение]
  На втором уроке начните с контрольных значений, затем выберите один из двух
  взглядов: графический способ для осмысления числа нулей или метод общих
  корней для сокращения вычислений. Самостоятельную реконструкцию используйте
  как проверку освоения, а задания на перенос --- для дифференциации.
\end{teacherbox}

\Needspace{7\baselineskip}
\begin{teacherbox}[Где важно не подсказывать]
  Дайте ученикам самостоятельно закончить: вторую ветвь модуля, объединение
  условий на \(a\), вторую подстановку в знаменатель и окончательную запись
  ответа. Именно эти места показывают, освоен ли метод, а не только прочитан
  образец.
\end{teacherbox}

\section{Ответы к рабочему листу}

\subsection{Правило и прогноз}

\begin{propertybox}
  \[
    \frac{A(x)}{B(x)}=0
    \quad\Longleftrightarrow\quad
    \begin{cases}
      A(x)=0,\\
      B(x)\ne0.
    \end{cases}
  \]
  После раскрытия модуля числитель задаётся двумя линейными выражениями,
  поэтому на каждой ветви возможно не более одного корня, а всего --- не
  более двух.
\end{propertybox}

\subsection{Исследование числителя}

При \(x\ge0\):
\[
  3x-2x=a+2
  \quad\Longrightarrow\quad
  x_1=a+2,
  \qquad a\ge-2.
\]

При \(x<0\):
\[
  -3x-2x=a+2
  \quad\Longrightarrow\quad
  x_2=-\frac{a+2}{5},
  \qquad a>-2.
\]

Уравнение числителя имеет два различных корня при \(a>-2\). При \(a=-2\) обе
формулы дают \(0\), но значение \(0\) не удовлетворяет условию \(x<0\),
поэтому остаётся один ноль.

\begin{conclusionbox}
  Числитель имеет ровно два различных нуля тогда и только тогда, когда
  \[
    a>-2.
  \]
\end{conclusionbox}

\subsection{Проверка знаменателя}

Для \(x_1=a+2\):
\begin{align*}
  D(x_1)
  &=(a+2)^2-2(a+2)-a\\
  &=a^2+a\\
  &=a(a+1).
\end{align*}
Поэтому исключаются \(a=-1\) и \(a=0\).

Для \(x_2=-\dfrac{a+2}{5}\):
\begin{align*}
  D(x_2)
  &=\frac{(a+2)^2}{25}+\frac{2(a+2)}5-a\\
  &=\frac{a^2-11a+24}{25}\\
  &=\frac{(a-3)(a-8)}{25}.
\end{align*}
Поэтому исключаются \(a=3\) и \(a=8\).

\begin{resultbox}
  \[
    \boxed{a\in(-2,+\infty)\setminus\{-1,0,3,8\}}.
  \]
\end{resultbox}

\parameternumberlinefigure

\subsection{Контрольные значения}

\begin{center}
  \renewcommand{\arraystretch}{1.42}
  \begin{tabularx}{\textwidth}{
    >{\centering\arraybackslash}p{0.12\textwidth}
    >{\centering\arraybackslash}p{0.24\textwidth}
    >{\raggedright\arraybackslash}p{0.42\textwidth}
    >{\centering\arraybackslash}X}
    \toprule
    \rowcolor{TableHeader}
    \(a\) & Нулей числителя & Знаменатель & Корней дроби \\
    \midrule
    \(-2\) & 1 & При \(x=0\): \(D(0)=2\ne0\) & 1 \\
    \addlinespace[0.2em]
    \(-1\) & 2 & Кандидат \(x_1=1\) запрещён & 1 \\
    \addlinespace[0.2em]
    \(1\) & 2 & Оба значения допустимы & 2 \\
    \bottomrule
  \end{tabularx}
\end{center}

\section{Графическое объяснение}

Положим \(b=a+2\). График функции
\[
  y=3|x|-2x=
  \begin{cases}
    -5x,&x<0,\\
    x,&x\ge0
  \end{cases}
\]
имеет вершину в начале координат и целиком лежит не ниже оси \(Ox\).

\parametergraphfigure

Отсюда:
\[
  b<0\Rightarrow0\text{ корней},\qquad
  b=0\Rightarrow1\text{ корень},\qquad
  b>0\Rightarrow2\text{ корня}.
\]
Так как \(b=a+2\), два нуля числителя возникают при \(a>-2\).

\begin{warningbox}
  График относится только к числителю. Он не содержит информации об области
  определения исходного рационального выражения, поэтому значения
  \(-1,0,3,8\) на нём не обнаруживаются.
\end{warningbox}

\section{Исключение общих корней}

Общий ноль числителя и знаменателя удовлетворяет системе
\[
  \begin{cases}
    |3x|-2x-2-a=0,\\
    x^2-2x-a=0.
  \end{cases}
\]
Выразив \(a\), получаем
\[
  a=3|x|-2x-2,
  \qquad
  a=x^2-2x,
\]
следовательно,
\[
  x^2=3|x|-2.
\]

При \(x\ge0\):
\[
  x^2-3x+2=0
  \quad\Longrightarrow\quad
  x=1,2
  \quad\Longrightarrow\quad
  a=-1,0.
\]

При \(x<0\):
\[
  x^2+3x+2=0
  \quad\Longrightarrow\quad
  x=-1,-2
  \quad\Longrightarrow\quad
  a=3,8.
\]

Таким способом сразу находятся четыре значения параметра, которые нужно
исключить из условия \(a>-2\).

\section{Типичные ошибки и реакция учителя}

\begin{center}
  \renewcommand{\arraystretch}{1.42}
  \begin{tabularx}{\textwidth}{
    >{\raggedright\arraybackslash}p{0.32\textwidth}
    >{\raggedright\arraybackslash}p{0.28\textwidth}
    >{\raggedright\arraybackslash}X}
    \toprule
    \rowcolor{TableHeader}
    Ошибка & Диагностика & Наводящий вопрос \\
    \midrule
    Сразу записано \(a>-2\)
      & Исследован только числитель
      & Что должно быть верно для знаменателя при найденном \(x\)? \\
    \addlinespace[0.2em]
    Для отрицательной ветви записано \(a\ge-2\)
      & Потеряно строгое условие \(x<0\)
      & Может ли \(x=0\) принадлежать отрицательной ветви? \\
    \addlinespace[0.2em]
    При \(a=-2\) получены два корня
      & Две одинаковые формулы приняты за разные корни
      & Какие реальные значения \(x\) получаются? \\
    \addlinespace[0.2em]
    Исключаются значения \(x\), а не \(a\)
      & Смешаны переменная и параметр
      & Что именно требуется найти в условии задачи? \\
    \addlinespace[0.2em]
    График объявлен полным решением
      & Не проверена область определения
      & Как на графике числителя увидеть нули знаменателя? \\
    \bottomrule
  \end{tabularx}
\end{center}

\section{Дифференциация}

\begin{teacherbox}[Базовый маршрут]
  Оставьте разделы 1--4 ученического листа и самостоятельную реконструкцию.
  Графическое объяснение проведите устно, а поиск общих корней не включайте.
\end{teacherbox}

\begin{teacherbox}[Сильная группа]
  После аналитического решения предложите сначала самостоятельно найти
  исключаемые параметры через систему общих корней. Затем сравните число
  вычислений в двух методах и обсудите, какой способ легче обобщать.
\end{teacherbox}

\begin{teacherbox}[Второй урок или кружок]
  Постройте полную классификацию числа корней и рассмотрите задачу, в которой
  параметр находится внутри модуля. Там изменяется положение точки излома,
  поэтому прежний графический шаблон уже не переносится механически.
\end{teacherbox}

\section{Ответы к заданиям на перенос}

\begin{propertybox}[Задание 1]
  Для уравнения
  \[
    \frac{|2x|-x-1-a}{x^2-x-a}=0
  \]
  числитель имеет два нуля при \(a>-1\):
  \[
    x_1=a+1,
    \qquad
    x_2=-\frac{a+1}{3}.
  \]
  Проверка знаменателя даёт исключения \(a=0\) и \(a=2\). Ответ:
  \[
    \boxed{a\in(-1,+\infty)\setminus\{0,2\}}.
  \]
\end{propertybox}

\begin{propertybox}[Задание 2]
  Условие \(a>-2\) не является ложным: оно правильно описывает число
  нулей числителя. Решение неполно, потому что не проверено, являются ли эти
  нули допустимыми для исходного уравнения.
\end{propertybox}

\clearpage
\section{Критерии освоения}

\begin{questionsbox}
  \small
  Ученик освоил метод, если без опоры на готовый текст:
  \par\vspace{0.35em}
  \begin{minipage}[t]{0.485\linewidth}
    \begin{itemize}[itemsep=0.18em,topsep=0.15em]
      \item формулирует условие нуля дроби;
      \item получает обе ветви и сохраняет ограничения;
      \item объясняет случай \(a=-2\);
    \end{itemize}
  \end{minipage}\hfill
  \begin{minipage}[t]{0.485\linewidth}
    \begin{itemize}[itemsep=0.18em,topsep=0.15em]
      \item подставляет оба найденных нуля числителя в знаменатель;
      \item различает промежуточное и итоговое условия;
      \item проверяет ответ на контрольном значении.
    \end{itemize}
  \end{minipage}
\end{questionsbox}


\subsection{Быстрая оценка самостоятельного решения}

\begin{center}
  \renewcommand{\arraystretch}{1.42}
  \begin{tabularx}{\textwidth}{
    >{\raggedright\arraybackslash}p{0.42\textwidth}
    >{\centering\arraybackslash}p{0.16\textwidth}
    >{\centering\arraybackslash}p{0.16\textwidth}
    >{\centering\arraybackslash}X}
    \toprule
    \rowcolor{TableHeader}
    Критерий & 2 балла & 1 балл & 0 баллов \\
    \midrule
    Две ветви модуля и их ограничения
      & Верно & Одна неточность & Не построены \\
    \addlinespace[0.2em]
    Условие двух различных нулей
      & Обосновано & Записано без границы & Неверно \\
    \addlinespace[0.2em]
    Проверка двух найденных нулей в знаменателе
      & Обе проверки & Одна проверка & Нет проверки \\
    \addlinespace[0.2em]
    Окончательный ответ
      & Полный & Есть техническая ошибка & Не соответствует решению \\
    \bottomrule
  \end{tabularx}
\end{center}

\begin{teacherbox}[Интерпретация]
  7--8 баллов: метод освоен; 5--6 баллов: требуется точечная коррекция;
  3--4 балла: ученик понимает отдельные шаги, но не удерживает архитектуру;
  0--2 балла: нужно вернуться к правилу нуля дроби и двум слоям задачи.
\end{teacherbox}

\subsection{Лист наблюдения за самостоятельной работой}

\begin{center}
  \small
  \renewcommand{\arraystretch}{1.55}
  \begin{tabularx}{\textwidth}{
    >{\raggedright\arraybackslash}p{0.27\textwidth}
    >{\centering\arraybackslash}p{0.12\textwidth}
    >{\centering\arraybackslash}p{0.15\textwidth}
    >{\centering\arraybackslash}p{0.12\textwidth}
    >{\raggedright\arraybackslash}X}
    \toprule
    \rowcolor{TableHeader}
    Ученик & Ветви & Проверка \(D(x)\) & Ответ & Комментарий \\
    \midrule
    \rule{0pt}{5mm} & & & & \\
    \midrule
    \rule{0pt}{5mm} & & & & \\
    \midrule
    \rule{0pt}{5mm} & & & & \\
    \midrule
    \rule{0pt}{5mm} & & & & \\
    \midrule
    \rule{0pt}{5mm} & & & & \\
    \midrule
    \rule{0pt}{5mm} & & & & \\
    \midrule
    \rule{0pt}{5mm} & & & & \\
    \midrule
    \rule{0pt}{5mm} & & & & \\
    \midrule
    \rule{0pt}{5mm} & & & & \\
    \midrule
    \rule{0pt}{5mm} & & & & \\
    \bottomrule
  \end{tabularx}
\end{center}

\begin{teacherbox}[Итоговая заметка]
  Что потребовало повторного объяснения? Какие ученики смогли перейти к
  методу общих корней? Какой шаг стоит изменить в следующей версии листа?
\end{teacherbox}
\workarea{28mm}
